\subsection{Sprint 3}
Durante o período estipulado para o desenvolvimento das tarefas da terceira sprint, estavam previstas as seguintes \ac{us}: \ac{us}01 (Cadastro de Atrações Turísticas), \ac{us}03 (Visualização de Atração Turística), \ac{us}05 (Tela de Home), \ac{us}08 (Visualização de Evento), além de resolver alguns débitos técnicos da sprint anterior, como popular o banco de dados, filtragem da lista de experiências, header da home, imagens e armazenamento ECII. Fiquei encarregada, junto de outros colegas, da tarefa de integração do envio do formulário para cadastro na \ac{aws}.

Infelizmente, devido às consequências contínuas da enchente que ainda indisponibilizavam a participação de diversos integrantes da equipe, a task de integração do envio do formulário para cadastro na \ac{aws} não pôde ser finalizada. Mesmo assim, a sprint foi entregue com suas User Stories concluídas e os débitos técnicos pendentes da sprint anterior foram resolvidos. Assim, busquei me fazer presente e disponível para ajudar em outras tarefas, como a criação do fluxo de cadastro de um evento no Figma.

Durante essa sprint, continuamos a enfrentar os efeitos da grande enchente que atingiu o Rio Grande do Sul, resultando na impossibilidade de diversos colegas contribuírem com a equipe. Isso causou o adiamento de algumas tarefas para a próxima sprint. Além disso, o início do meu período de estágio trouxe desafios adicionais, como a redução do tempo disponível para me dedicar ao projeto. A conciliação entre faculdade e trabalho se mostrou difícil, fazendo com que eu sentisse que não consegui contribuir tanto quanto desejava.

Apesar dos desafios, aprendemos, como equipe, que devemos manter a determinação, fazendo o possível e dando o nosso melhor com os recursos disponíveis. Esta lição de resiliência e adaptabilidade foi essencial para a continuidade do projeto, mesmo em tempos de adversidade.
Depois da entrega, buscamos finalizar os débitos técnicos desta sprint para garantir a continuidade do desenvolvimento das tasks previstas para as próximas semanas. Além disso, busquei fazer um planejamento de rotina que permita uma maior dedicação e contribuição de minha parte para com o projeto nessa reta final.